% Bibliografia: "onde comprar ou baixar" a partir dos campos de apostila.dbx
% (url, acesso, urlgratis, edicaopt, urlpt, doi, eprint). Convenção no
% cabeçalho de cada bib/*.bib.

% url que é só o DOI ou o arXiv fica no .bib para o link check (lychee),
% mas não é impresso: o DOI e o eprint já aparecem com link.
\DeclareSourcemap{\maps[datatype=bibtex]{
  \map{\step[fieldsource=url, match=\regexp{\Ahttps://doi\.org/}, final]
       \step[fieldset=url, null]\step[fieldset=urldate, null]}
  \map{\step[fieldsource=eprinttype, match=\regexp{\Aarxiv\Z}, final]
       \step[fieldset=url, null]\step[fieldset=urldate, null]}
}}

% Texto do campo acesso (sempre depois de ponto na bibliografia).
\newcommand*{\acessotexto}[1]{%
  \ifstrequal{#1}{compra}{À venda na editora}{%
  \ifstrequal{#1}{aberto}{Acesso aberto}{%
  \ifstrequal{#1}{pago}{Acesso pago}{%
  \ifstrequal{#1}{esgotado}{Esgotado}{#1}}}}}
\DeclareFieldFormat{acesso}{\acessotexto{#1}}
\DeclareFieldFormat{urlpt}{\url{#1}}
\DeclareFieldFormat{urlgratis}{\url{#1}}

% No fim de cada registro da bibliografia: acesso, edição PT, leitura grátis.
\renewbibmacro*{finentry}{%
  \iffieldequalstr{langid}{brazilian}{\setunit{\addspace}\printtext{\pt}}{}%
  \iffieldundef{acesso}{}{\setunit{\addperiod\space}%
    \printtext{\mkbibemph{\printfield{acesso}}}}%
  \iffieldundef{edicaopt}{}{\setunit{\addperiod\space}%
    \printtext{Edição PT: \printfield{edicaopt}}%
    \iffieldundef{urlpt}{\setunit{\addspace}\printtext{(sem loja da editora)}}
      {\setunit{\addcomma\space}\printfield{urlpt}}}%
  \iffieldundef{urlgratis}{}{\setunit{\addperiod\space}%
    \printtext{Leitura gratuita: \printfield{urlgratis}}}%
  \finentry}

% ---- Listas "Onde achar conteúdo" (\printbiblist{ondeachar}) ----
% Uma linha por registro, colunas geradas dos campos; nada de URL à mão.
\DeclareBiblistFilter{ondeachar}{\filter[type=field, filter=title]}
\newtoggle{ondecompacto}% papers, vídeos, artigos, cursos: só Obra e Link
\newcommand*{\ondecel}[2]{\parbox[t]{#1\linewidth}{\raggedright#2}}
\newcommand*{\ondelinhas}[4]{% Obra | Comprar | Edição PT | Download gratuito
  \iftoggle{ondecompacto}%
    {\ondecel{.74}{#1}\hfill\ondecel{.24}{#4}}%
    {\ondecel{.37}{#1}\hfill\ondecel{.17}{#2}\hfill
     \ondecel{.24}{#3}\hfill\ondecel{.18}{#4}}}
\defbibenvironment{ondeachar}
  {\list{}{\setlength{\leftmargin}{0pt}\setlength{\itemindent}{0pt}%
     \setlength{\itemsep}{4pt}\setlength{\parsep}{0pt}}\small}
  {\endlist}
  {\item\relax}
% Cabeçalho das colunas, impresso antes de cada lista.
\newcommand{\ondecabecalho}{\par\noindent\begingroup\small\bfseries
  \iftoggle{ondecompacto}{\ondelinhas{Obra}{}{}{Link}}%
    {\ondelinhas{Obra}{Comprar}{Edição PT}{Download gratuito}}%
  \par\endgroup\nobreak\vspace{2pt}\hrule\nobreak\vspace{3pt}}

% Prefixos de link montados aqui; "https:" separado do resto para que o
% link check (lychee) não trate o modelo como endereço.
\edef\ondedoibase{\detokenize{https:}//doi.org/}
\edef\ondearxivbase{\detokenize{https:}//arxiv.org/abs/}
\newbibmacro{onde:obra}{%
  \mkbibemph{\printfield[noformat]{title}}%
  \iffieldequalstr{langid}{brazilian}{\addspace\pt}{}%
  \ifnameundef{labelname}{}{\newline\printnames{labelname}}%
  \iffieldundef{year}{}{\addcomma\space\printfield{year}}}
\newbibmacro{onde:comprar}{%
  \iffieldequalstr{acesso}{compra}
    {\href{\thefield{url}}{\iflistundef{publisher}{editora}{\printlist{publisher}}}}
  {\iffieldequalstr{acesso}{esgotado}
    {esgotado\iffieldundef{url}{}{ (\href{\thefield{url}}{catálogo})}}
    {---}}}
\newbibmacro{onde:pt}{%
  \iffieldequalstr{langid}{brazilian}{original em português}
  {\iffieldundef{edicaopt}{---}
    {\iffieldundef{urlpt}{\printfield{edicaopt} (sem loja da editora)}
      {\href{\thefield{urlpt}}{\printfield{edicaopt}}}}}}
% Link aberto: urlgratis; senão url, arXiv ou DOI quando acesso é aberto
% (ou ausente, como em vídeos e documentação); acesso pago: o DOI, mais o
% urlgratis quando há versão aberta em relatório técnico (outro documento).
\newbibmacro{onde:gratis}{%
  \iffieldequalstr{acesso}{pago}
    {\iffieldundef{urlgratis}{}%
       {\href{\thefield{urlgratis}}{relatório técnico grátis}; }%
     \href{\ondedoibase\thefield{doi}}{acesso pago (DOI)}}
  {\iffieldundef{urlgratis}
  {\ifboolexpr{test {\iffieldundef{acesso}} or test {\iffieldequalstr{acesso}{aberto}}}
    {\iffieldundef{url}
      {\iffieldequalstr{eprinttype}{arxiv}
        {\href{\ondearxivbase\thefield{eprint}}{arXiv}}
        {\iffieldundef{doi}{---}{\href{\ondedoibase\thefield{doi}}{DOI}}}}
      {\href{\thefield{url}}{\ifkeyword{video}{assistir}{abrir}}}}
    {---}}
  {\href{\thefield{urlgratis}}{ler grátis}}}}
\DeclareBibliographyDriver{ondeachar}{%
  \ondelinhas{\usebibmacro{onde:obra}}{\usebibmacro{onde:comprar}}%
    {\usebibmacro{onde:pt}}{\usebibmacro{onde:gratis}}}
% Uso: \ondeachar{livro}  ou  \ondeachar*{video}  (* = compacto)
\NewDocumentCommand{\ondeachar}{s m}{%
  \IfBooleanTF{#1}{\toggletrue{ondecompacto}}{\togglefalse{ondecompacto}}%
  \ondecabecalho
  \printbiblist[keyword=#2, heading=none]{ondeachar}}

% URLs longas quebram em qualquer caractere na bibliografia (evita overfull).
\appto{\bibsetup}{\sloppy}
\setcounter{biburlnumpenalty}{9000}
\setcounter{biburllcpenalty}{7000}
\setcounter{biburlucpenalty}{8000}
